% GENERATED FILE. Edit canonical lessons, then run npm run generate:books.
% canonical-source-hash: fnv1a64:023df525688f92d2
% canonical-lessons: ZH-C158-jiaoxing, ZH-C158-shuizhao, ZH-C158-zuomeng, ZH-C158-qilai, ZH-C158-xiguan

\chapter{Waking Up}
\label{ch:zh-waking-up}
\hlchaptermodality{158}

\begin{chapteropening}
\textbf{By the end of this chapter:} \emph{I can say who wakes me, when I fall asleep, what I dream and what I am used to.}

Complete the last lesson of chapter 158: I can say who wakes me, when I fall asleep, what I dream and what I am used to.
\end{chapteropening}

\section[jiàoxǐng]{jiàoxǐng — to wake someone up}

\label{lesson:ZH-C158-jiaoxing}



\begin{quote}
Before the new one: say the Mandarin for news, then the Mandarin for to reply.
\end{quote}

\subsection*{You'll want to know: jiàoxǐng}
\emph{jiàoxǐng} — \textquotedblleft{}to wake someone up\textquotedblright{}.

\begin{quote}
\emph{māma měitiān qī diǎn jiàoxǐng wǒ} — Mum wakes me at seven every day.
\end{quote}

Its characters come later; for now it is a word for the ear.

\begin{cousinweb}
\emph{jiàoxǐng} is built from \zh{叫} 'to call' and \emph{xǐng} 'to wake', both inherited from Middle Chinese, \emph{kewH} and \emph{sengX} (Guangyun, in Baxter's transcription). The dictionary glosses the second as sobering up after wine: to call someone awake.
\end{cousinweb}

\begin{grammarlens}[title={What you've built}]
One more word for saying what I do.
\end{grammarlens}

\subsection*{Your turn}
\begin{itemize}
\raggedright
  \item \emph{Say it:} \emph{jiàoxǐng}
  \item \emph{Say it:} \emph{jiàoxǐng}, once more
  \item \emph{Say it:} \emph{māma měitiān qī diǎn jiàoxǐng wǒ} — Mum wakes me at seven every day.
  \item \emph{Recall:} say the Mandarin for news, then the Mandarin for to reply, then say \emph{jiàoxǐng} again
\end{itemize}

\begin{tcolorbox}[breakable,colback=teal!4,colframe=teal!35!black,title={Before you move on}]
What does jiàoxǐng mean? (\textquotedblleft{}To wake someone up\textquotedblright{}.) Say it once more.
\end{tcolorbox}
\section[shuìzháo]{shuìzháo — to fall asleep}

\label{lesson:ZH-C158-shuizhao}



\begin{quote}
Before the new one: say the Mandarin for to reply, then the Mandarin for to wake someone up.
\end{quote}

\subsection*{You'll want to know: shuìzháo}
\emph{shuìzháo} — \textquotedblleft{}to fall asleep\textquotedblright{}.

\begin{quote}
\emph{wǒ zuótiān wǎnshang shí'èr diǎn shuìzháo le} — Last night I fell asleep at twelve.
\end{quote}

Its characters come later; for now it is a word for the ear.

\begin{cousinweb}
\emph{shuìzháo} is built from \emph{shuì} 'to sleep' and \emph{zháo}, both inherited from Middle Chinese, \emph{dzyweH} and \emph{drjak} (Guangyun, in Baxter's transcription). The dictionary glosses the second 'to attach to': sleep takes hold.
\end{cousinweb}

\begin{grammarlens}[title={What you've built}]
One more word for saying what I do.
\end{grammarlens}

\subsection*{Your turn}
\begin{itemize}
\raggedright
  \item \emph{Say it:} \emph{shuìzháo}
  \item \emph{Say it:} \emph{shuìzháo}, once more
  \item \emph{Say it:} \emph{wǒ zuótiān wǎnshang shí'èr diǎn shuìzháo le} — Last night I fell asleep at twelve.
  \item \emph{Recall:} say the Mandarin for to reply, then the Mandarin for to wake someone up, then say \emph{shuìzháo} again
\end{itemize}

\begin{tcolorbox}[breakable,colback=teal!4,colframe=teal!35!black,title={Before you move on}]
What does shuìzháo mean? (\textquotedblleft{}To fall asleep\textquotedblright{}.) Say it once more.
\end{tcolorbox}
\section[zuòmèng]{zuòmèng — to dream; to have a dream}

\label{lesson:ZH-C158-zuomeng}



\begin{quote}
Before the new one: say the Mandarin for to wake someone up, then the Mandarin for to fall asleep.
\end{quote}

\subsection*{You'll want to know: zuòmèng}
\emph{zuòmèng} — \textquotedblleft{}to dream; to have a dream\textquotedblright{}.

\begin{quote}
\emph{wǒ chángcháng zuòmèng} — I dream a lot.
\end{quote}

Its characters come later; for now it is a word for the ear.

\begin{cousinweb}
\emph{zuòmèng} is built from \zh{做} 'to do' and \emph{mèng} 'a dream'. The second is inherited from Middle Chinese \emph{mjuwngH} (Guangyun, in Baxter's transcription), glossed as the spirit wandering in sleep; \zh{做} has no form in that dictionary, so the history of the first part is not settled.
\end{cousinweb}

\begin{grammarlens}[title={What you've built}]
One more word for saying what I do.
\end{grammarlens}

\subsection*{Your turn}
\begin{itemize}
\raggedright
  \item \emph{Say it:} \emph{zuòmèng}
  \item \emph{Say it:} \emph{zuòmèng}, once more
  \item \emph{Say it:} \emph{wǒ chángcháng zuòmèng} — I dream a lot.
  \item \emph{Recall:} say the Mandarin for to wake someone up, then the Mandarin for to fall asleep, then say \emph{zuòmèng} again
\end{itemize}

\begin{tcolorbox}[breakable,colback=teal!4,colframe=teal!35!black,title={Before you move on}]
What does zuòmèng mean? (\textquotedblleft{}To dream; to have a dream\textquotedblright{}.) Say it once more.
\end{tcolorbox}
\section[\texorpdfstring{\zh{起来} (qǐlai)}{起来 (qǐlai)}]{\zh{起来} (qǐlai) — to get up; to stand up}

\label{lesson:ZH-C158-qilai}



\begin{quote}
Before the new one: say the Mandarin for to fall asleep, then the Mandarin for to have a dream.
\end{quote}

\subsection*{You'll want to know: \zh{起来}}
\textbf{\zh{起来}} — \emph{qǐlai} — \textquotedblleft{}to get up; to stand up\textquotedblright{}.

\begin{quote}
\textbf{\zh{七点了你快起来}} \emph{qī diǎn le nǐ kuài qǐlai} — It's seven o'clock: up you get, quick.
\end{quote}

\begin{cousinweb}
\textbf{\zh{起来}} (\emph{qǐlai}) is built from \zh{起} 'to rise' and \zh{来} 'to come', both inherited from Middle Chinese, \emph{khiX} and \emph{loj} (Guangyun, in Baxter's transcription), which glosses the first 'to rise' and 'to stand'. In the word the second loses its tone.
\end{cousinweb}

\begin{grammarlens}[title={What you've built}]
One more word for saying what I do.
\end{grammarlens}

\subsection*{Your turn}
\begin{itemize}
\raggedright
  \item \emph{Say it:} \emph{qǐlai}
  \item \emph{Say it:} \emph{qǐlai}, once more
  \item \emph{Say it:} \emph{qī diǎn le nǐ kuài qǐlai} — It's seven o'clock: up you get, quick.
  \item \emph{Recall:} say the Mandarin for to fall asleep, then the Mandarin for to have a dream, then say \emph{qǐlai} again
\end{itemize}

\begin{tcolorbox}[breakable,colback=teal!4,colframe=teal!35!black,title={Before you move on}]
What does \zh{起来} mean? (\textquotedblleft{}To get up; to stand up\textquotedblright{}.) Say it once more.
\end{tcolorbox}
\section[xíguàn]{xíguàn — to be used to; a habit}

\label{lesson:ZH-C158-xiguan}



\begin{quote}
Before the new one: say the Mandarin for to have a dream, then the Mandarin for to stand up.
\end{quote}

\subsection*{You'll want to know: xíguàn}
\emph{xíguàn} — \textquotedblleft{}to be used to; a habit\textquotedblright{}.

\begin{quote}
\emph{wǒ hái méi xíguàn liù diǎn qǐchuáng} — I'm still not used to getting up at six.
\end{quote}

Its characters come later; for now it is a word for the ear.

\begin{cousinweb}
\emph{xíguàn} is built from \emph{xí} and \emph{guàn}, both inherited from Middle Chinese, \emph{zip} and \emph{kwaenH} (Guangyun, in Baxter's transcription). The dictionary glosses the first 'to learn' and explains the second with the first: practised.
\end{cousinweb}

\begin{grammarlens}[title={What you've built}]
One more word for saying what I do.
\end{grammarlens}

\subsection*{Your turn}
\begin{itemize}
\raggedright
  \item \emph{Say it:} \emph{xíguàn}
  \item \emph{Say it:} \emph{xíguàn}, once more
  \item \emph{Say it:} \emph{wǒ hái méi xíguàn liù diǎn qǐchuáng} — I'm still not used to getting up at six.
  \item \emph{Recall:} say the Mandarin for to have a dream, then the Mandarin for to stand up, then say \emph{xíguàn} again
\end{itemize}

\begin{tcolorbox}[breakable,colback=teal!4,colframe=teal!35!black,title={Before you move on}]
What does xíguàn mean? (\textquotedblleft{}To be used to; a habit\textquotedblright{}.) Say it once more.
\end{tcolorbox}
